\section{Test-time search, and the statistic that decides whether it helps}
\label{sec:search}

Every prior FishBot chooses its ask by maximising a fitted score over the legal set. The obvious
next step, and the one the project owner asked for, is to reason forward: to evaluate a candidate by
what happens after it rather than by a static function of the position it is played from. This
section builds that and measures it. It is the one mechanism in this study that beats a static
rule --- and getting there required correcting a statistic whose cost, \vsixCurseGap{} points, is
the most transferable result in the paper.

\subsection{Why determinization here is not double-dummy search}
\label{sec:search-notpimc}

The v0.4 study eliminated perfect-information Monte Carlo on structural grounds and this study
re-confirms them: under perfect information the team on turn takes every half-suit not dealt outright
to its opponents, so a double-dummy evaluator scores every legal action alike and the average over
determinizations cannot discriminate. Fish is a game about who knows what; a clairvoyant evaluator
prices none of it.

The search built here determinizes the \emph{deal} --- which is the entire hidden state because every card movement in this game is public --- but does not make anybody clairvoyant. Each of the six players in the
continuation is reconstructed \emph{at its own information set}: the public deduction state, which
is common knowledge and is therefore one object rather than six, refined by the hand that
determinization gives that seat. They miss asks, emit certificates and misdeclare in the proportions
the real game does. This is the cooperative single-agent search of Lerer et al.\
\cite{lerer2020sparta}, with the belief taken from an exactly enumerable posterior rather than a learnt
one.

The reconstruction is checked against ground truth rather than argued. Over \vsixReconStates{} states and \vsixReconChecks{} card checks, the reconstructed
information set differs from the seat's real one on \vsixReconMaskPct\% of cards and is \emph{wider}
in \vsixReconWider{} of \vsixReconWider{} cases (narrower in \vsixReconNarrower{}), replicated at
\vsixReconBanks{} seed banks: it is a strict under-approximation of
what the seat knows, so no simulated player is ever credited with knowledge it does not have.

\subsection{The search}
\label{sec:search-alg}

At an ask decision the actor
(i) builds the exact capacity dynamic program on its own deduction state and draws $D$ deals from
it, rejecting draws that violate an ask-legality certificate, so accepted draws are exact samples of
the policy-agnostic posterior, and re-weighting them by the policy prior as an importance weight so
that the sampler stays exact and the prior stays auditable;
(ii) takes the $K$ leading candidates under the blueprint score, extended through every candidate
that ties the $K$-th, because a plain top-$K$ would cut an exact tie arbitrarily and so reproduce the
very defect the search is being asked about;
(iii) seats six blueprint agents at their reconstructed information sets, applies the candidate, and
plays the continuation out;
(iv) and repeats every candidate on the \emph{same} determinizations, so the comparison is paired.

Two design choices are forced by measurement rather than taste. The blueprint used inside the
rollout is the deployed policy with its two-ply refinement disabled and its Sinkhorn iteration count
reduced; Table~\ref{tab:rollout} gives the cost/fidelity frontier, which the corpus had not swept.
The cheapest blueprint available is $62\times$ cheaper per event than the deployed policy and
$38$ points weaker, and a search whose continuation model is $38$ points weak is searching a
different game.

\begin{table}[t]\centering
\caption{Rollout-blueprint fidelity. Win rate against \vfive{} and single-thread throughput.
Reducing the Sinkhorn iteration count, which the corpus had not tried, is the cheap axis: it buys
most of the speed of the independent-marginal blueprint at a small fraction of its strength cost.}
\label{tab:rollout}
\begin{tabular}{lrrr}
\toprule
Rollout blueprint & Win rate vs \vfive{} & 95\% CI & Games/s \\
\midrule
\texttt{v05:belief=indep,topk=0} & 12.00 & [10.33, 13.75] & 996.2 \\
\texttt{v05:topk=0,souter=1,sinner=1} & 46.08 & [43.33, 48.83] & 146.8 \\
\texttt{v05:topk=0,souter=1,sinner=3} & 46.75 & [44.00, 49.50] & 114.9 \\
\texttt{v05:topk=0,souter=2,sinner=4} & 51.42 & [48.67, 54.17] & 78.7 \\
\texttt{v05:topk=0} & 47.75 & [44.92, 50.50] & 30.7 \\
\bottomrule
\end{tabular}

\end{table}

\subsection{The result: the optimizer's curse is worth \vsixCurseGap{} points}
\label{sec:search-curse}

The first working build of the search chose the candidate with the highest mean rollout return and
scored \vsixSearchUnguarded\% against \vfive{}. The control that isolates the cause is the same
code, the same rollout blueprint, the same budget, the same seed bank and the same sample size, with
the blueprint's preference made to dominate: it scores \vsixSearchControl\%. The machinery is
correct and the \emph{statistic} was wrong, and the gap between those two rows ---
\vsixCurseGap{} points --- is the price of the argmax.

One caveat about that control, because it is the row the whole comparison rests on. The blueprint's
preference dominates only where the blueprint \emph{has} one: inside its own tie group the added
term is identically zero and the deviation penalty is waived, so the control still lets the search
choose among tied candidates. It scores at parity, which is a second, independent confirmation of
\S\ref{sec:diag-irreducible}: searching inside an exchangeable tie group changes nothing.

One determinization's return is the final half-suit differential, whose standard deviation is about
two and a half half-suits, while genuine differences between the leading candidates are an order of
magnitude smaller. An argmax over $D$ such means therefore selects on residual noise and lands on a
candidate the blueprint had ranked below the top. Replacing the argmax with a paired
lower-confidence-bound rule --- deviate from the blueprint's own choice only when the improvement
over it, on the same determinizations, clears $\kappa$ standard errors of the paired difference ---
recovers the whole deficit monotonically in $\kappa$ (Table~\ref{tab:search}).

\begin{table}[t]\centering
\caption{The deviation-threshold ladder, six rotations, held-out bank. The first row is the control
that forces the blueprint to decide. The range across the ladder is \vsixSearchSpread{} points.}
\label{tab:search}
\begin{tabular}{lrrr}
\toprule
Configuration & Win rate vs \vfive{} & 95\% CI & $n$ \\
\midrule
\texttt{s1=1,det=8,cand=6,blend=1000000} & 49.31 & [45.83, 52.78] & 720 \\
\texttt{s1=1,det=8,cand=6,kappa=0} & 13.61 & [10.97, 16.39] & 720 \\
\texttt{s1=1,det=12,cand=4,kappa=0} & 23.33 & [20.42, 26.39] & 720 \\
\texttt{s1=1,det=12,cand=4,kappa=1} & 42.78 & [39.31, 46.25] & 720 \\
\texttt{s1=1,det=12,cand=4,kappa=2.5} & 53.75 & [50.42, 57.22] & 720 \\
\texttt{s1=1,det=12,cand=4,kappa=4} & 53.89 & [50.42, 57.36] & 720 \\
\texttt{s1=1,det=16,cand=6,kappa=2,maxq=26} & 48.89 & [45.14, 52.64] & 720 \\
\texttt{s1=1,det=16,cand=6,kappa=2,maxq=26,deadsearch=2} & 50.28 & [46.67, 53.89] & 720 \\
\bottomrule
\end{tabular}

\end{table}

This is a property of determinized evaluation in this game rather than of this implementation, and
it is stated as the section's main contribution. Any search over a high-variance terminal return
that selects by an unguarded argmax will pay it, and the price here --- \vsixSearchSpread{} points
across the ladder --- is larger than every mechanism the three prior studies fitted, combined.

\subsection{Where the advantage lives}
\label{sec:search-falsifier}

Guarded, the search stops losing and starts winning: at $\kappa = 2.5$ it takes
\vsixSearchPooled\% against \vfive{} pooled over \vsixSearchPooledBanks{} banks and
\vsixSearchPooledN{} games. The remaining
question is whether that win is the kind a working search produces, and the deviation rate answers
it --- but only once the rate is decomposed. The falsifier the literature supplies is explicit: a healthy test-time search differs from its
blueprint on a small percentage of decisions, and one that differs on tens of percent is mis-scaled
rather than smart. At $\kappa = 2.5$ the search moves the action on \vsixSearchDevTwoFive\% of the decisions it
searches (Table~\ref{tab:searchdev}), and that rate does not fall below about $31\%$ however far
$\kappa$ is raised. Read naively this is an order of magnitude above the band the literature
associates with a search that refines a good policy, and it was so read in an earlier draft. It is
the wrong reading, and the reason is \S\ref{sec:diag-irreducible}.

The deviation penalty is waived inside the blueprint's own tie group, because there the blueprint
expresses no preference. That group is more than half of all decisions, so the search always
re-chooses inside it, and those re-choices are neutral by construction --- the candidates are
exchangeable. The floor of $31\%$ is that churn, not signal. Applying the same penalty inside the
tie group collapses the rate to \vsixSearchDevGuardedTwoFive\% at $\kappa = 2.5$ and
\vsixSearchDevGuardedFour\% at $\kappa = 4$.

That decomposition also bears on where the win sits. Applying the penalty inside the tie group --- so that the search deviates only on the
\vsixSearchDevGuardedTwoFive\% of decisions where it has evidence against the blueprint's stated
preference --- takes the win rate to \vsixSearchTieGuardedWin\%
[\vsixSearchTieGuardedLo, \vsixSearchTieGuardedHi]: the advantage disappears. Forbidding deviation
\emph{outside} the tie group instead, so the search only ever re-chooses among candidates the
blueprint cannot separate, takes it to \vsixSearchTieOnlyTwelve\%
[\vsixSearchTieOnlyTwelveLo, \vsixSearchTieOnlyTwelveHi] --- better than the unrestricted search.

\textbf{What is established, and what is not.} The pooled effect is established twice over. On
\vfive{}'s parameter vector the guarded search takes \vsixSearchPooled\% against \vfive{} over
\vsixSearchPooledN{} games and \vsixSearchPooledBanks{} banks, none below
\vsixSearchPooledMin\%. On \vsix{}'s own vector --- so that the search is being asked to improve
the best static policy this project has, not the previous one --- it takes \vsixSearchLift\%
against \vsix{} over \vsixSearchLiftN{} games, with \vsixSearchLiftAbove{} of
\vsixSearchLiftCells{} cells above parity and the worst at \vsixSearchLiftMin\%. That second
result is what makes \vsixsearch{} a configuration worth naming. So are the two negative
controls of \S\ref{sec:diag-irreducible}: randomising the tie group is worth exactly zero and one
sampled deal is worth nothing, so what the ensemble is reading is joint information and not the
mere fact of varying a deterministic choice.

What is \emph{not} established is where inside the search the effect sits. Guarding the tie group
gives \vsixSearchTieGuardedWin\%; restricting the search to the tie group gives
\vsixTieSearchBankList\% across \vsixTieSearchBanks{} banks; the unrestricted search gives
\vsixSearchPooled\%. At 720 games a cell those intervals overlap each other and the unrestricted
figure, and an earlier draft of this section read a single favourable bank as an attribution. It is
not one. Resolving it needs roughly an order of magnitude more games at a configuration that costs
\vsixSearchGps{} games per second on one thread, against \vsixThroughputSix{} for the deployed
policy across all threads --- which is the reason it is not resolved here.

\begin{table}[t]\centering
\caption{Search deviation rate against the deviation threshold. A healthy search moves a small
percentage of decisions; this one either moves a third of them or none.}
\label{tab:searchdev}
\begin{tabular}{lrrr}
\toprule
$\kappa$ & Decisions searched (\%) & Moved the action (\%) & Games/s \\
\midrule
0 & 98.58 & 65.67 & 0.173 \\
1 & 98.51 & 41.51 & 0.165 \\
2.5 & 98.51 & 33.49 & 0.144 \\
4 & 98.32 & 31.40 & 0.144 \\
6 & 98.32 & 31.59 & 0.144 \\
\bottomrule
\end{tabular}

\end{table}



\subsection{The one move class that exists only under search}
\label{sec:search-deadask}

M1 restricts the candidate set to asks the actor cannot prove will miss. That fixed v0.4's
two-question cycle, and it deleted the move class every multi-step human tactic is built out of:
blackballing, choosing which opponent receives the turn, and paying for a signal with a safe ask are
all a guaranteed miss at a \emph{chosen} seat. Such an ask exists at $79.2\%$ of \vfive{}'s
decisions and at two or more distinct chosen opponents at $50.1\%$.

Unbanning it wholesale is measured and rejected. With dead asks re-admitted and the ownership
features gated by hit probability so that the $p=0$ incentive behind v0.4's cycle cannot fire, the
configuration scores \vsixDeadWin\% [\vsixDeadWinLo, \vsixDeadWinHi] against \vfive{} --- higher ---
and fails the commit gate outright: \vsixDeadPathDead\% of its asks are provably dead, its longest
dead run is \vsixDeadPathLongest{}, and \vsixDeadPathLimit\% of its games are killed by the action
limit. This is the trap \S\ref{sec:protocol} exists to prevent, and it is the same trap the v0.5
ablation table contains.

Rationing it does not help either, and the reason is instructive: \emph{the linear score cannot
price a deliberate miss}. With the four ownership features zeroed on a dead candidate, the admission
margin was swept at four settings and the output was bit-identical at every one, because the hit
probability term is zero on such an ask by definition and every remaining term of the score is a
penalty. Its value is the value of the position the donated turn produces, which no static feature
of the ask computes.

The search can price it, and this is the clearest case in the study of a move class that exists only
under search: the deliberate miss is offered to the rollouts as a candidate, one per distinct target,
under an anti-repeat guard that is provably free. A dead $(c,t)$ pair stays dead unless the card
publicly moves to $t$, and that movement is observable, so forbidding the unmoved repeat forbids
nothing that could have become live. This is exactly the distinction the blanket repetition guard of
the v0.5 study missed, and why that guard cost six points: it also forbade \emph{live} repeats, which
are frequently the best ask on the board once a card has moved.

With the deliberate miss in the candidate set the endgame-restricted search scores
\vsixSearchDead\%. It is the best search configuration measured and it is still not separated from
the blueprint.

\subsection{What ships}
\label{sec:search-ships}

The search ships \textbf{off}, and the reason is cost and unresolved scope rather than a null
result. On the ladder's best setting it is ahead of \vfive{} with an interval that excludes parity,
but it is three orders of magnitude slower than the deployed policy and it moves a third of all
decisions --- so what it is is not a refinement of the
shipped policy but a distinct, far more expensive one whose evaluation this study could not complete
at the resolution \S\ref{sec:protocol} demands of everything else. It is retained, switchable and
fully instrumented. \S\ref{sec:limitations} states what would have to be measured before it could
ship on.
